\section{Literature Review}%
\label{sec:literature_review}

This review examines twenty papers spanning three interconnected research
threads: (1)~model-based testing (MBT) and its practical/industrial
adoption, (2)~statistical model checking (SMC) as a scalable alternative
to exhaustive probabilistic verification, and (3)~probabilistic and
uncertainty-aware extensions that bridge testing and formal verification.
Each paper is discussed in turn, covering its context, contribution,
methodology, and relevance to research on probabilistic model-based
testing under uncertainty.


%  Papers 1--5: Model-Based Testing


\subsection{Test Model Coverage Analysis Under Uncertainty}
\label{sec:lr_prasetya}

Prasetya and Klomp~\cite{prasetya2021coverage} directly target the
problem of measuring test coverage when the underlying test model is
non-deterministic a situation that arises either because the model is
an abstraction of the real system or because the system under test is
itself non-deterministic.  In such settings a single test case may
trigger different execution paths depending on internal decisions of the
software, which makes exact coverage computation impossible.  The
authors' solution is to let developers annotate each model transition
with an estimated probability, enabling a probabilistic model-checking
algorithm to compute \emph{probabilistic coverage}.  The paper's key
extension over prior model-checking approaches is that it moves beyond
simple \emph{reachability}-style coverage queries (which can be answered
by standard probabilistic model checkers) toward \textbf{aggregate
coverage} goals e.g., ``80\% of states covered'' and
\textbf{$k$-wise coverage}, which is known from the testing literature
to substantially increase fault-detection potential.  Because aggregate
and $k$-wise goals cannot be expressed in LTL/CTL and are not answerable
by conventional model checkers, the authors present a new, efficient
labelling algorithm to compute these probabilistic aggregate/$k$-wise
coverage measures directly.  This paper is highly relevant to any
research studying probabilistic MBT under uncertainty, as it provides
one of the few rigorous formal treatments of \emph{coverage measurement
itself} becoming a probabilistic quantity when models are
non-deterministic a foundational concern for any downstream
uncertainty-aware testing framework.

\subsection{Practitioners' Best Practices for MBT Adoption}
\label{sec:lr_alegroth}

Al\'{e}groth et al.~\cite{alegroth2022practitioners} conducted an
industrial, interview-based empirical study addressing a persistent
puzzle in the MBT literature: despite decades of academic research,
industrial adoption of graphical-model-based MBT remains sparse.  The
authors conducted semi-structured interviews with 17~international MBT
experts across different industrial roles, then synthesised the results
through semantic equivalence analysis, later verified by additional
practitioners.  The study yields 13~synthesised conclusions and
23~concrete best-practice guidelines covering the full lifecycle of
adoption, use, and (where appropriate) abandonment of graphical MBT.
Importantly, the paper's conclusions span not just technical dimensions
(model design, tool choice, maintainability) but organisational and
process factors mindset, knowledge, mandate, and resource
allocation that the authors argue are just as decisive in determining
whether MBT succeeds in practice.  This paper is valuable to the
literature not for its formal/mathematical contribution but as a
counterweight: it grounds more theoretical work on probabilistic
coverage and uncertainty-aware testing in the practical reality that most
industrial MBT failures are organisational, not algorithmic.

\subsection{Transforming Model-Based Tests into Code}
\label{sec:lr_ferrari}

Ferrari et al.~\cite{ferrari2023transforming} present a systematic
literature review (SLR) investigating a question central to the
practical value of MBT: does coverage achieved at the abstract model
level actually translate into meaningful coverage of the underlying
source code once abstract tests are transformed into concrete, executable
tests?  Using a snowballing methodology starting from three seed papers,
the authors expanded their search to 30~primary studies.  The review
characterises how test sets generated from models are mapped onto code,
catalogues the transformation techniques used across the literature, and
critically examines the (surprisingly thin) empirical evidence connecting
model coverage criteria to code coverage outcomes.  As an SLR, its main
contribution to the field is a structured map of the model-to-code
transformation landscape together with an explicit identification of
gaps most notably the lack of large-scale empirical validation that
model-level testing effort reliably produces code-level assurance.  This
is directly relevant background for any argument that model coverage
analysis (as in~\cite{prasetya2021coverage}) needs empirical grounding in
code-level outcomes, and it complements industrial accounts of similar
transformation pipelines in practice~\cite{garousi2021mbt_practice}.

\subsection{Generative Model-Based Testing on Decision-Making Policies}
\label{sec:lr_li}

Li et al.~\cite{li2023generative} tackle the modern problem of testing
decision-making policies that solve Markov decision processes
(MDPs) for example, DNN-based policies deployed in autonomous driving
or robotics.  Because such policies operate over deep-neural-network-based,
effectively infinite state spaces, exhaustive or even conventional
model-based test generation is impractical.  The authors propose a
generative testing framework built around a \textbf{diffusion-model-based
test case generator}, chosen for its ability to adapt flexibly to
different search spaces while producing valid, in-distribution test
cases.  A key algorithmic contribution is a
\textbf{termination-state novelty-based guidance mechanism}, which steers
the generative process toward diversifying agent behaviour at episode
termination, thereby increasing the likelihood of discovering diverse and
influential failure-triggering scenarios.  The framework is evaluated
across five benchmarks, including autonomous driving and aircraft
collision avoidance, and is shown to surface more diverse and impactful
failures than prior state-of-the-art baselines.  This paper represents
the modern generative-AI evolution of model-based test generation replacing
hand-crafted model traversal strategies with learned generative models and
is squarely relevant to any survey that wants to show how classical MBT
ideas (models of behaviour, coverage-directed exploration) have been
re-instantiated using generative deep learning for MDP-based systems.

\subsection{Coverage Measurement in MBT of Web Applications}
\label{sec:lr_garousi_coverage}

Garousi et al.~\cite{garousi2024coverage} present an applied,
tool-oriented paper arising from a large-scale industrial
web-application-testing context.  The authors identify a concrete
practical gap: existing coverage tools typically report only one type of
coverage (either code coverage or requirements/test coverage), whereas
practitioners running large MBT test suites need to integrate front-end
code coverage, back-end code coverage, and requirements coverage into a
single ``live'' view as tests execute.  Unable to find any suitable
off-the-shelf toolset, the authors built an open-source tool,
\textbf{MBTCover}, purpose-built for MBT contexts, which reports code,
requirements, and model coverage simultaneously and in real time during
test execution.  The paper presents the tool's architecture and the
authors' experience deploying it across multiple large test-automation
projects.  As a direct, applied companion to~\cite{prasetya2021coverage}
(which is theoretical) and~\cite{garousi2021mbt_practice} (a broader MBT
experience report from the same research group), this paper demonstrates
that the practical challenges of coverage measurement in MBT are as much
about tooling and integration as about the underlying probabilistic or
combinatorial theory.


%  Papers 6--7: Bridging MBT and SMC


\subsection{PASTA: Proactive Adaptation via Statistical Model Checking}
\label{sec:lr_pasta}

Shin et al.~\cite{shin2021pasta} address proactive adaptation in
self-adaptive systems (SAS) adaptation triggered by \emph{predicting}
future environmental changes rather than reacting after the fact.
Because predictions are inherently uncertain, adaptation consequences
must be verified before being executed, and prior work relied on
probabilistic model checking (PMC) for this verification.  The authors
identify two limitations of PMC-based verification in this setting: the
state-explosion problem for complex SAS models, and restrictive support
for particular modelling languages.  PASTA replaces PMC with
\textbf{statistical model checking (SMC)}, generating statistically
sufficient samples to verify the effects of adaptation tactics far faster
than exhaustive PMC approaches.  The paper contributes not just an
algorithm but a reference architecture and an open-source implementation
skeleton, evaluated on two real self-adaptive systems using actual
operational data, alongside a comparative analysis of the relative
advantages/disadvantages of PMC versus SMC for proactive adaptation.
This paper is a strong example of SMC being adopted specifically to
overcome the scalability limitations of exhaustive probabilistic
verification precisely the trade-off that motivates much of the broader
SMC literature~\cite{agha2018survey, legay2010smc, larsen2014smc}.

\subsection{Uncertainty-Aware Exploration in Model-Based Testing}
\label{sec:lr_camilli}

Camilli et al.~\cite{camilli2021uncertainty} propose what is among the
most directly relevant contributions to the theme of probabilistic MBT
under uncertainty, since it explicitly frames uncertainty as a first-class
testing concern rather than an incidental complication.  The authors
propose novel model-based exploration strategies for generating test
cases that specifically target the \emph{uncertain} components of a
system under test, using \textbf{Markov Decision Processes} as the
underlying modelling formalism.  Testers explicitly attach
\emph{beliefs} (probability distributions) to transition probabilities
to represent their confidence (or lack thereof) in the model's fidelity
to the real system.  The structural properties of the model, together
with the uncertainty specification, drive test-case generation, and
\textbf{Bayesian inference} is used to update the initial beliefs as
evidence accumulates from testing.  The paper introduces several concrete
test-selection strategies (Flat/random, History-based, Distance-based,
Frequency-based) and empirically evaluates them on synthetic systems of
varying structural complexity, finding that the best strategy depends on
system complexity and the overall uncertainty level.  This paper together
with a related follow-on piece envisioning online,
reinforcement-learning-based uncertainty testing represents one of the
clearest bridges in the literature between classical MBT and formal
treatment of epistemic uncertainty via Bayesian updating.


%  Papers 8--17: Statistical Model Checking


\subsection{Sound Statistical Model Checking for Probabilities and Expected Rewards}
\label{sec:lr_budde}

Budde et al.~\cite{budde2025sound} provide a rigorous, foundational
contribution to the theory underpinning SMC tools.  The authors observe
that many existing and even newly developed SMC tools are either
\textbf{unsound} (their statistical guarantees are not actually valid,
often because confidence intervals are computed via central-limit-theorem
approximations that do not hold in the finite-sample regime) or
\textbf{inefficient} (relying on overly conservative bounds such as the
Okamoto bound, requiring far more samples than necessary).  The paper
provides a comprehensive survey of existing SMC tools' correctness and
of sound statistical estimation methods from the literature, applied to
the problem of estimating both probabilities and \textbf{expected
rewards}.  For expected rewards specifically, the authors show how to
bound the path-reward distribution so that sound statistical methods
designed for bounded distributions can be applied; they recommend the
\textbf{Dvoretzky--Kiefer--Wolfowitz (DKW) inequality}, previously
unused in SMC, and formally prove that even reachability rewards can, in
theory, be bounded introducing the concept of \emph{limit-PAC
procedures} as a practical solution when exact bounds are unavailable.
These methods are implemented in the ``modes'' SMC tool and
experimentally validated.  This paper is essential reading for
understanding the current (2025) state-of-the-art in statistically
rigorous SMC, and it directly informs any research that wants to make
quantitative uncertainty claims about test/verification outcomes with
genuine statistical guarantees.

\subsection{Statistical Model Checking: The 2024 Edition}
\label{sec:lr_kanav}

Kanav et al.~\cite{kanav2025smc2024} present a short
introductory/overview note accompanying a dedicated SMC session at
AISoLA~2024.  It briefly reintroduces the core ideas of statistical
model checking and surveys the field's trajectory effectively an
updated, contemporary companion to the earlier ``past, present, and
future'' retrospectives~\cite{larsen2014smc}.  While brief, the note
situates recent SMC developments (e.g., PAC statistical model checking
of mean payoff, decision-tree-based controller representation via
\emph{dtcontrol}, and applications to security-critical code analysis)
within the longer arc of the field, and signals which open problems the
community considers current priorities.  Its main value to a literature
review is as a fast-moving pulse-check on the field as of 2024,
complementing the more technically dense contemporary
papers~\cite{budde2025sound, soltani2023spare} with a broader,
session-introduction-level perspective.

\subsection{PRISM: Statistical Model Checking in Practice}
\label{sec:lr_prism}

The official PRISM documentation~\cite{prism_smc} explains how PRISM's
discrete-event simulator can approximate property values via sampling
rather than exact numerical solution a technique particularly valuable
for very large models where exhaustive model checking is infeasible,
since simulation proceeds directly from the PRISM language description
without constructing the full underlying probabilistic model.  The
documentation details the two default statistical methods available:
\emph{Confidence Interval} estimation for quantitative properties, and
the \emph{Sequential Probability Ratio Test} (SPRT) for bounded
properties.  It also specifies the configurable parameters (sample count,
confidence level, interval width) and the practical restrictions on which
properties and modelling-language features SMC in PRISM supports
(currently limited to P/R operators, no LTL-style path properties or
filters).  While not an original research contribution, this
documentation is an important practical reference underpinning much of
the applied SMC literature reviewed here.

\subsection{Rigorous Processor Evaluation with Statistical Model Checking}
\label{sec:lr_processors}

The paper on rigorous evaluation of computer
processors~\cite{smc_processors2023} introduces SMC to an entirely new
application domain: \textbf{computer architecture and processor
evaluation}.  The authors observe that experiments with real computer
processors necessarily exhibit variability, and that prior work injects
randomness into simulators to mimic this variability, generating
distributions of benchmark results across multiple runs.  However,
existing evaluation practice typically applies standard statistical
techniques that implicitly (and often incorrectly) assume these result
distributions are Gaussian.  To enable rigorous evaluation for arbitrary,
unknown result distributions, the authors adapt statistical model
checking to processor performance analysis, framing the ``SMC for
Processor Analysis'' (SPA) methodology.  SMC's statistical guarantees
allow architecture researchers to determine, with a specified confidence
level, whether a system satisfies a performance property for at least a
given fraction of executions, without assuming any particular underlying
distribution.  This paper is a good illustration of SMC's methodological
portability: the same statistical machinery developed for verifying
stochastic system models is directly repurposed for empirical systems
evaluation.

\subsection{Optimal Spare Management via Statistical Model Checking}
\label{sec:lr_soltani}

Soltani et al.~\cite{soltani2023spare} apply statistical model checking
to a concrete industrial reliability-engineering problem: determining the
optimal number of spare parts to stock in order to balance the twin goals
of system reliability and cost.  The authors combine \textbf{fault tree
analysis} with SMC by modelling spare-part management as a
\textbf{stochastic priced timed game automaton (SPTGA)}, then use
\textbf{Uppaal Stratego} to find the spare-part count that minimises
total costs from downtime and purchasing, while the resulting SPTGA model
can also be queried for other metrics such as expected availability.  The
technique is applied to the emergency shutdown system of a research
nuclear reactor a rare-event setting in which the relevant failure
probabilities are extremely low, requiring the authors to adjust Uppaal
Stratego's settings to obtain statistically reliable estimates despite
the rarity of the events of interest.  The reported result an optimal
spare configuration achieving 99.96\% expected availability, computed in
minutes rather than the days required by exhaustive
alternatives exemplifies SMC's core value proposition (tractable
analysis of otherwise intractable stochastic models) in a genuinely
safety-critical, real-world application.

\subsection{PRG4CNN: Probabilistic Robustness Guarantees for CNNs}
\label{sec:lr_prg4cnn}

Liu and Fang~\cite{liu2025prg4cnn} extend probabilistic model checking
into the domain of neural-network robustness verification.  Motivated by
the well-known fragility of convolutional neural networks (CNNs) to
small input perturbations, and noting that existing robustness
verification methods focus predominantly on \emph{local} robustness
(which is computationally expensive) and cannot automatically
\emph{repair} a network once a robustness violation is found, the authors
propose \textbf{PRG4CNN} described as the first automated and complete
framework for guaranteeing \emph{probabilistic} robustness of CNNs.  The
framework operates in four steps: (1)~model the CNN as a
\textbf{Markov Decision Process} via model learning, (2)~specify the
desired probabilistic robustness property using \textbf{PCTL}
(Probabilistic Computational Tree Logic), (3)~verify this property using
a probabilistic model checker, and (4)~if the property does not hold,
perform \textbf{probabilistic robustness repair} via
counterexample-guided sensitivity analysis.  This paper is relevant to
the broader review as an example of probabilistic model checking being
extended from classical reactive/stochastic-system verification toward
machine learning model assurance an increasingly important adjacent
application area for the model-checking and MBT communities.

\subsection{Statistical Model Checking Meets Property-Based Testing}
\label{sec:lr_aichernig}

Aichernig and Schumi~\cite{aichernig2017smc} present a key
methodological bridge between the testing community and the
formal-verification-oriented SMC community one of the most direct
precedents for combining statistical model checking with mainstream
software testing practice.  The authors note that SMC scales well to
large stochastic models and is relatively simple to implement precisely
because it avoids the state-space-explosion problem inherent to
exhaustive model checking, by instead simulating finitely many executions
and applying hypothesis testing to infer whether the observed samples
support or refute a given property.  Their contribution is to show how
SMC can be integrated into a \textbf{property-based testing}
framework specifically \textbf{FsCheck} for C\# yielding a flexible
combined testing/simulation environment in which programmers define both
models and properties in an ordinary programming language, without
needing an external, specialised modelling language.  This design has two
key advantages: it eliminates the need for a separate modelling
formalism, and it allows both stochastic \emph{models} and actual
\emph{implementations} to be checked using the same property-based
machinery.  This work is an important conceptual precursor for any
framework that wants to unify probabilistic testing with statistical
formal verification~\cite{camilli2021uncertainty}.

\subsection{A Survey of Statistical Model Checking}
\label{sec:lr_agha}

Agha and Palmskog~\cite{agha2018survey} present one of the most
comprehensive and widely cited surveys of the SMC field, covering
algorithms, techniques, and tools for statistical model checking of
stochastic systems whose quantitative properties are typically specified
in probabilistic temporal logics.  The survey systematically covers both
\textbf{hypothesis-testing-based} approaches (e.g., the sequential
probability ratio test, following Younes and Sen et~al.) and
\textbf{estimation-based} approaches (including Bayesian estimation
methods, such as those of Zuliani et~al., which bound the probability of
an incorrect answer using a prior distribution and posterior estimation),
as well as variance-reduction techniques like importance sampling with
the cross-entropy method for analysing rare-event properties.  The
authors explicitly emphasise the trade-offs between precision and
scalability that define the practical usability of different SMC
algorithms.  As a survey published in 2018, this paper offers the most
thorough single reference point for understanding the algorithmic
landscape that underlies nearly every applied SMC paper in this
review~\cite{shin2021pasta, budde2025sound, smc_processors2023,
soltani2023spare, aichernig2017smc}.

\subsection{Statistical Model Checking: An Overview}
\label{sec:lr_legay}

Legay et al.~\cite{legay2010smc} present an earlier, foundational
tutorial-style paper that is one of the field's most frequently cited
introductions to statistical model checking.  The authors contrast the
traditional \textbf{numerical approach} to model checking stochastic
systems which iteratively computes or approximates the exact measure of
executions satisfying a temporal-logic subformula with the
\textbf{statistical/simulation-based approach}: simulating the system for
a finite number of executions and applying hypothesis testing to infer
whether the observed samples provide statistical evidence for or against
the property of interest.  The paper's stated purpose is to survey this
statistical approach and articulate its principal advantages:
\emph{efficiency} (avoiding exhaustive state-space construction),
\emph{uniformity} (applicability across many modelling formalisms, since
only a simulator is required), and \emph{simplicity} (conceptually
straightforward compared to symbolic or numerical probabilistic model
checking algorithms).  While largely superseded in technical depth by
later, more comprehensive surveys~\cite{agha2018survey} and later
foundational advances in soundness~\cite{budde2025sound}, this 2010
paper remains historically important as one of the field-defining
tutorial references establishing SMC as a distinct research area separate
from exhaustive probabilistic model checking.

\subsection{Statistical Model Checking: Past, Present, and Future}
\label{sec:lr_larsen}

Larsen and Legay~\cite{larsen2014smc} provide an invited track
introduction accompanying a dedicated ``Statistical Model Checking:
Past, Present, and Future'' session at ISoLA~2014.  The paper provides
historical framing for the SMC field, tracing its development from early
formal-methods work through to (as of 2014) contemporary applications
and tools.  The authors characterise SMC as fundamentally a
\textbf{compromise between formal verification and testing}: system
executions are monitored (as in testing) until a statistical algorithm
can produce a rigorous estimate of the property of interest (echoing the
rigour traditionally associated with exhaustive formal verification
methods, which by contrast can ensure full correctness for all possible
scenarios but at far greater computational cost).  Read together
with~\cite{kanav2025smc2024, agha2018survey, legay2010smc}, this
paper despite being over a decade old at the time of this
review completes a useful chronological arc documenting how the SMC
field's foundational concerns (efficiency, soundness, scope of
applicability) evolved from 2010 through to the 2023--2025 wave of
papers reviewed above.


%  Papers 18--20: Probabilistic MBT Foundations


\subsection{Model-Based Testing of Probabilistic Systems}
\label{sec:lr_gerhold}

Gerhold and Stoelinga~\cite{gerhold2016probabilistic} present a
foundational paper for probabilistic model-based testing, directly
connecting classical \textbf{ioco-theory} (input/output conformance
testing) with \textbf{statistical hypothesis testing}.  The authors
present an executable MBT framework for black-box probabilistic systems
that also exhibit non-determinism, providing algorithms to automatically
generate, execute, and evaluate test cases from a probabilistic
requirements model.  Functionally, the ioco-style algorithms handle
correctness of the traditional (non-probabilistic) input/output
behaviour; statistically, $\chi^{2}$~hypothesis tests and
distribution-fitting methods assess whether the \emph{frequencies}
observed during repeated test execution correspond to the
\emph{probabilities} specified in the requirements model.  A key
technical challenge the authors solve is that non-determinism in the
underlying probabilistic input/output transition systems (pIOTS)
prevents a direct application of the $\chi^{2}$~test; instead, the
authors formulate a non-linear optimisation problem to find the ``best
resolution'' of the non-determinism that could explain the observed test
data, enabling the $\chi^{2}$~test to then be applied meaningfully.  The
paper's central theoretical results are the \textbf{soundness} (every
test case receives the mathematically correct verdict) and
\textbf{completeness} (the framework can, in principle, detect any
probabilistic deviation from the specification with arbitrary precision)
of the resulting \textbf{probabilistic input/output conformance (pioco)}
relation.  This paper is arguably the most rigorous formal foundation
among all twenty papers reviewed here for probabilistic model-based
testing specifically, and is essential background
for~\cite{prasetya2021coverage}'s coverage-analysis extensions
and~\cite{camilli2021uncertainty}'s uncertainty-aware exploration
strategies.

\subsection{MBT in Practice: Web Applications Domain}
\label{sec:lr_garousi_practice}

Garousi et al.~\cite{garousi2021mbt_practice} present an industrial
experience report documenting the deployment of model-based testing at a
large software testing company (Testinium A.\c{S}.) to elevate the
maturity of its test-automation practices for large-scale web and mobile
applications.  Based on an ``action research'' methodology, the authors
selected the open-source tool \textbf{GraphWalker} from among available
open-source/commercial MBT tools and pragmatically applied MBT for
end-to-end test automation.  The reported benefits are both tangible
(improved test coverage measured as number of distinct paths tested, and
improved real-fault detection effectiveness) and intangible (improved
test-design discipline across the organisation).  The paper's central
argument echoed by~\cite{alegroth2022practitioners}'s findings on
adoption barriers is that the practical value of any MBT technique is
inseparable from its usability by working test engineers under real
resource constraints (time, effort, existing skillsets, management
priorities), and that heavyweight, academically sophisticated MBT
approaches often fail in enterprise contexts like web/mobile testing
precisely because they neglect this.  This paper's authorship overlaps
substantially with~\cite{garousi2024coverage} (the MBTCover tool paper),
and together the two form a coherent, practice-grounded narrative of a
real industrial MBT adoption journey.

\subsection{Approximate Planning and Verification for Large MDPs}
\label{sec:lr_lassaigne}

Lassaigne and Peyronnet~\cite{lassaigne2012approximate} address the
planning and verification problems for very large probabilistic
systems specifically Markov decision processes from a computational
complexity perspective.  The authors' central concern is designing
\textbf{efficient approximation methods} for two related problems:
(1)~computing a near-optimal policy for the planning problem in
discounted MDPs, and (2)~computing the satisfaction probabilities of
properties of interest (such as reachability or safety) over the Markov
chain that results from restricting the MDP to that near-optimal policy.
Two distinct approximation approaches are presented: the first is based
on \textbf{sparse sampling}, whose complexity is notably independent of
the size of the state space and requires only a probabilistic generator
of the MDP rather than full access to its transition structure making
it applicable even when the state space is too large to represent
explicitly; the second uses a variant of the \textbf{multiplicative
weights update algorithm}.  The paper provides a complete complexity
analysis of the sparse-sampling approach, parameterised primarily by the
desired approximation quality.  Because MDPs are the common underlying
formalism connecting several other papers in this
review \cite{li2023generative}'s generative testing of MDP-based
decision policies, \cite{camilli2021uncertainty}'s uncertainty-aware MBT
over MDPs, and~\cite{liu2025prg4cnn}'s MDP-based CNN robustness
verification this paper's complexity-theoretic treatment of MDP
planning/verification approximation provides useful theoretical grounding
for why sampling-based (rather than exhaustive) approaches are necessary
once MDP-based models scale beyond small, hand-crafted examples.


%  Synthesis


\subsection{Synthesis}
\label{sec:lr_synthesis}

Taken together, these twenty papers trace three converging lines of
research.  First, the \textbf{model-based testing}
tradition~\cite{prasetya2021coverage, alegroth2022practitioners,
ferrari2023transforming, li2023generative, garousi2024coverage,
gerhold2016probabilistic, garousi2021mbt_practice} has matured from
academic proposals toward serious grappling with practical adoption
barriers, coverage-measurement rigour, and the model-to-code
transformation gap.  Second, the \textbf{statistical model checking}
tradition~\cite{budde2025sound, kanav2025smc2024, prism_smc,
smc_processors2023, soltani2023spare, agha2018survey, legay2010smc,
larsen2014smc} has evolved from foundational tutorials (2010--2014)
through a comprehensive survey (2018) to a current wave (2023--2025) of
papers establishing genuine statistical soundness and extending SMC into
entirely new domains computer architecture, spare-parts logistics,
quantum calibration that share nothing with SMC's original
formal-verification roots except the underlying statistical machinery.
Third, a smaller but conceptually pivotal set of
papers~\cite{shin2021pasta, camilli2021uncertainty, liu2025prg4cnn,
aichernig2017smc, lassaigne2012approximate} explicitly \textbf{bridges}
testing and statistical/probabilistic verification, using MDPs, PCTL, and
hypothesis testing as the common mathematical vocabulary.  Collectively,
they support the view that treating uncertainty as a first-class,
quantifiable concern rather than an obstacle to be abstracted away is
now a mainstream, multi-disciplinary concern spanning software testing,
formal methods, computer architecture, and machine-learning assurance.
